\begin{figure}[h]
    \centering
    \includegraphics[width=\linewidth]{figures/audio/MovieGen_Audio_Extension.pdf}
    \caption{\textbf{\mvga{} extension diagram.} A user provides a video (\eg, 58s), and audio caption for each video chunk (\eg, 20s). Starting from the second chunk, the model takes not only the video chunk and the caption, but also a segment from the previously generated audio (\eg, the last 5s) in order to generate a new chunk that is coherent with the previous one.
    }
    \label{fig:audio_sys}
\end{figure}

In order to generate soundtracks for variable duration of videos, we build
a single model that can perform both \textit{audio generation} given a video,
and \textit{audio extension} given a video with partially generated audio.
We aim to generate up to 30 seconds of audio in a single shot,
and allow the model to utilize extension to generate audio of arbitrary lengths.
\cref{fig:audio_sys} illustrates the process for long-form video generation.

We enable audio extension by training the model to perform masked audio prediction,
where the model predicts the audio target given the whole video and its surrounding audio.
The surrounding audio can be empty (\ie, audio generation),
before or after the target audio (\ie, audio extension in either direction),
or around the target (\ie, audio infilling).
Audio infilling is useful for fixing small segments that contains artifacts or unwanted sound effects.

Lastly, for sound design purposes, users would often want to specify what and how acoustic events should be added to the video,
such as deciding what on-screen sounds to emphasize, what off-screen sounds to add, whether there is background music, and what style to generate for the music.
To provide users more control, we enable text prompting.
